\section{The Recursive Spawning Protocol}
\label{sec:rs-protocol}

The Recursive Spawning Protocol (RSP) is the structural mechanism that enforces an immutable, auditable, and architecturally verifiable delegation chain from every active agent to a human principal. RSP is the second named pillar of the \bxthree{} Framework, operating across all three layers by design: the Purpose Layer provides the validation gate that authorizes each spawn event, the Bounds Engine enforces depth constraints as structural invariants, and the Fact Layer establishes the immutable parent pointer and maintains the tamper-evident forensic ledger. The protocol converts the delegation chain from a forensic reconstruction attempted after the fact into a runtime-verifiable artifact --- one that is immutable, auditable, and enforced at every spawn event with no bypass paths, no advisory phases, and no override mechanisms available to any machine actor under any operational circumstance.

The protocol executes through five mandatory structural phases, each mapping to a specific \bxthree{} layer mechanism. Each phase must complete successfully before the next begins. A failure at any phase terminates the spawn attempt. There is no partial compliance, no deferred enforcement, and no skip-on-success.

% --------------------------------------------------------------------------
\subsection{Phase 1: Purpose Layer Validation Gate}
\label{sec:rs-phase1-purpose}

Every spawn event begins with a Purpose Layer validation. This gate is not a policy convention --- it is a structural constraint enforced by the Fact Layer pre-commit hook, which rejects any provisioning request that does not have a matching Purpose Layer warrant recorded in the forensic ledger before the spawn executes. The Purpose Layer is the exclusive authorization origin for all chain growth in the RSP architecture.

\medskip
\noindent\textbf{Validation conditions.} A parent node $n_i$ that intends to spawn a child $n_{i+1}$ must submit a spawn request to the Purpose Layer declaring: (P1) the proposed child operating scope $R(n_{i+1})$, (P2) the specific operational condition that necessitates spawning, and (P3) a demonstration that the condition cannot be resolved within $R(n_i)$. The Purpose Layer validates three mandatory conditions before issuing a warrant:

\begin{enumerate}
    \item[(P1)] \textbf{Scope refinement.} $R(n_{i+1})$ must be a strict subset of $R(n_i)$:
    \begin{equation}
    R(n_{i+1}) \subset R(n_i).
    \label{eq:rs-scope-refinement}
    \end{equation}
    Lateral scope stability --- where $R(n_{i+1}) = R(n_i)$ --- is not authorized. The chain grows through scope \emph{narrowing}, not sideways expansion. Equality at any link constitutes a drift condition and invalidates the warrant. The subset relationship is not a policy convention --- it is a structural invariant enforced at every spawn event by the Fact Layer pre-commit hook.

    \item[(P2)] \textbf{Operational necessity.} The parent must declare, in the Fact Layer authorization ledger, the precise condition requiring child spawning. The declaration must identify why the condition cannot be resolved within $R(n_i)$ and why the child is the minimum necessary decomposition of the current task. Vague justifications do not satisfy this condition.

    \item[(P3)] \textbf{Minimum scope proportionality.} The child scope must be the narrowest scope sufficient to address the declared condition. The Purpose Layer does not authorize child scopes broader than operationally necessary. Scope beyond minimum necessary constitutes a proportionality violation and blocks warrant issuance.
\end{enumerate}

\medskip
\noindent\textbf{Warrant issuance.} When all three conditions are satisfied, the Purpose Layer issues a spawn warrant $w_i$ --- a cryptographically signed record written atomically to the Fact Layer authorization ledger:

\begin{equation}
w_i = \bigl\langle
\text{parent}=n_i,\;
\text{child}=n_{i+1},;
R(n_{i+1}),;
\text{justification},;
\text{timestamp},;
\sigma_{\text{PL}}
\bigr\rangle .
\label{eq:spawn-warrant}
\end{equation}

The warrant is the sole authorized precondition for spawn. No machine actor may initiate a spawn event without a matching warrant in the Fact Layer ledger --- the Fact Layer pre-commit hook rejects any provisioning request that lacks a valid warrant before the operation reaches the ledger write stage. This makes the Purpose Layer's exclusive terminus role a structural guarantee, not a configuration option.

\medskip
\noindent\textbf{Human anchor establishment.} At root provisioning --- when a human principal $h$ first authorizes the agent node $n_0$ --- the Purpose Layer records the human anchor $h(n_0) = h$ in the Fact Layer authorization ledger. The anchor is non-collapsing: for all nodes $n_j$ in any chain descended from $n_0$, $h(n_j) = h(n_0) = h$. The anchor does not change as the chain deepens. No machine actor can serve as a chain root.

% --------------------------------------------------------------------------
\subsection{Phase 2: Bounds Engine Depth Check}
\label{sec:rs-phase2-bounds}

When a Purpose Layer warrant is present in the ledger, the Bounds Engine performs the depth check as a mandatory structural enforcement step before the Fact Layer provisioning write is executed. The Bounds Engine is the \bxthree{} layer responsible for tracking and enforcing the depth of any active delegation chain against the architectural maximum $D_{\max}$.

\medskip
\noindent\textbf{Depth tracking.} The depth of a node $n$, denoted $\text{depth}(n)$, is defined as the number of edges in the delegation chain from $n$ to the root human anchor $n_0$, where $\text{depth}(n_0) = 0$. For any node $n_i$ with parent $n_{i-1}$, the depth is:
\begin{equation}
\text{depth}(n_i) = \text{depth}(n_{i-1}) + 1.
\label{eq:depth-definition}
\end{equation}

The Bounds Engine maintains the depth counter for every active node in the system. The counter is written atomically to the Fact Layer forensic ledger at each spawn event as part of the node's sealed state block.

\medskip
\noindent\textbf{Depth constraint enforcement.} The Bounds Engine enforces the depth bound as a structural invariant at every spawn event:
\begin{equation}
\text{depth}(n_i) + 1 \leq D_{\max}.
\label{eq:depth-bound}
\end{equation}
If Equation~\ref{eq:depth-bound} is violated, the Bounds Engine sets the node state to \texttt{ESCALATING} and transmits a mandatory escalation event to the human anchor via the Bailout Protocol trigger pathway. The spawn is refused. No override, no emergency bypass, and no discretionary authority to exceed $D_{\max}$ regardless of operational urgency exists within the Bounds Engine. The $D_{\text{ESCALATE}}$ threshold provides advance warning at $D_{\max} - 1$, requiring human confirmation before the depth limit is reached --- not only after it is breached.

\medskip
\noindent\textbf{Depth limit rationale.} $D_{\max}$ is set at system initialization based on the operational consequence domain of the deployment. High-stakes environments --- financial orchestration, clinical decision support, critical infrastructure --- require lower $D_{\max}$ values. The value is recorded in the Fact Layer authorization ledger and is immutable. The depth bound is not a suggestion; it is a structural invariant enforced at every spawn step by both the Bounds Engine and the Fact Layer pre-commit hook.

% --------------------------------------------------------------------------
\subsection{Phase 3: Fact Layer --- Immutable Pointer Creation and Forensic Ledger Recording}
\label{sec:rs-phase3-fact-layer}

With the Purpose Layer warrant verified and the Bounds Engine depth check passed, the Fact Layer executes the atomic provisioning write that creates the immutable parent pointer and records the spawn event in the forensic ledger.

\medskip
\noindent\textbf{Atomic provisioning write.} The Fact Layer atomically executes five operations as a single indivisible transaction:

\begin{enumerate}
    \item[(A1)] Create node record $n_{i+1}$ in the forensic ledger.
    \item[(A2)] Write $\alpha(n_{i+1}) = n_i$ into $n_{i+1}$'s sealed memory envelope.
    \item[(A3)] Set $n_{i+1}$'s authorized scope to $R(n_{i+1})$ as declared in the Purpose Layer warrant.
    \item[(A4)] Seal the memory envelope: encrypt with a Fact Layer-only symmetric key and compute an HMAC-SHA-256 signature over the ciphertext.
    \item[(A5)] Append the complete spawn event as a single ledger entry $\ell_j$ to the forensic ledger $\mathcal{L}$.
\end{enumerate}

The atomicity guarantee is critical: there is no temporal window in which $\alpha(n_{i+1})$ is established without a matching Purpose Layer warrant, or in which a warrant exists without a corresponding parent pointer. The write is all-or-nothing --- if any step fails, all five are rolled back and the spawn attempt is logged as a failed event.

\medskip
\noindent\textbf{Immutability of the parent pointer.} Once $\alpha(n_{i+1}) = n_i$ is written, it is structurally permanent. The Fact Layer write protocol defines exactly one valid write operation for $\alpha$: the initial provisioning write. The protocol defines \emph{no valid write operations} for modifying $\alpha(n)$ after provisioning. Any operation --- from any source, with any justification --- that attempts to overwrite, redirect, or delete $\alpha(n)$ after provisioning is rejected at the protocol level before execution reaches the ledger. The immutability is protocol-enforced, not access-control-enforced. In access-control models, immutability can be circumvented by an actor who acquires sufficient privileges. In the Fact Layer write protocol, the modification operation is structurally impossible regardless of privilege elevation, credential compromise, or software vulnerability. There is no administrative override, no emergency bypass, and no privileged actor capable of altering the chain after provisioning.

\medskip
\noindent\textbf{Forensic ledger recording.} The Fact Layer records every authorized spawn event as an append-only ledger entry in the forensic ledger $\mathcal{L}$. Each entry has the structure:

\begin{equation}
\ell_j = \bigl\langle
\text{index}=j,;
\text{node}=n_i,;
\text{event-type},;
\text{payload},;
\text{timestamp},;
\text{attribution},;
\text{hash}(\ell_{j-1})
\bigr\rangle .
\label{eq:ledger-entry}
\end{equation}

The hash-chaining ($\text{hash}(\ell_{j-1})$) establishes continuity: no entry can be inserted, deleted, reordered, or altered after the fact without breaking chain continuity --- a break detectable by any observer with ledger read access. The ledger is the authoritative source of truth for all authorized chain state. It is not a selectively maintained event log --- it is a first-class data structure whose integrity is enforced by the Fact Layer write protocol and verifiable at any runtime point.

The ledger records every authorized spawn event, every state transition, every Purpose Layer directive received, and every unresolved exception. The complete spawn topology --- including the warrant, the parent pointer, and the child configuration --- is committed atomically as a single indivisible ledger entry.

\medskip
\noindent\textbf{Pre-commit hook enforcement.} The Fact Layer pre-commit hook verifies two structural invariants before any spawn is recorded:

\begin{enumerate}
    \item[(FC1)] $R(n_{i+1}) \subseteq R(n_i)$ --- enforced as a structural invariant, not a policy convention.
    \item[(FC2)] $\text{depth}(n_{i+1}) \leq D_{\max}$ --- the depth constraint, re-verified at the Fact Layer even though it was checked at the Bounds Engine.
\end{enumerate}

If either condition fails, the Fact Layer rejects the spawn and logs the rejection as a ledger entry. The pre-commit hook is the structural mechanism that enforces scope refinement and depth constraints at every spawn event, providing a second-level enforcement gate independent of the Bounds Engine.

% --------------------------------------------------------------------------
\subsection{Phase 4: Child Activation with Immutable Parent}
\label{sec:rs-phase4-child-activation}

With the immutable parent pointer established and the spawn event recorded in the forensic ledger, the child node $n_{i+1}$ is activated within its provisioned worksheet --- a containerized logic set delivered over-the-air (OTA) from the parent node. The child is activated with three immutable properties:

\begin{enumerate}
    \item \textbf{Immutable parent pointer.} $\alpha(n_{i+1}) = n_i$ is a Fact Layer property intrinsic to $n_{i+1}$, established at provisioning and stored in its sealed memory envelope. The child carries its parent pointer as an immutable attribute of itself, not as a mutable external reference held by the parent.

    \item \textbf{Authorized scope.} $R(n_{i+1})$ is a strict subset of $R(n_i)$, as validated by the Purpose Layer and enforced by the Fact Layer pre-commit hook. The child operates within a narrower scope than its parent, preserving the intent-refinement property of the chain.

    \item \textbf{Sealed memory envelope.} The parent pointer and authorized scope are encrypted and HMAC-signed in $n_{i+1}$'s sealed memory envelope. The envelope is decrypted only for authorized read operations by the Fact Layer; the plaintext is never exposed to the Bounds Engine or the child node's own code.
\end{enumerate}

The child is isolated from its parent in the sense that the parent's operational state (active, paused, terminated, decommissioned) does not affect $n_{i+1}$'s ability to reference its parent through $\alpha(n_{i+1})$. A parent that has terminated or crashed does not affect the child's chain-traversal capability. The child can always traverse outward to its complete lineage regardless of the state of any ancestor.

The child is also isolated from sibling nodes: no lateral links exist in a valid RSP chain, as each node has exactly one parent and one child scope relationship. The chains are nested rather than intersecting. This isolation is what makes RSP chains resistant to cross-node contamination and ensures that accountability attribution is always unambiguous.

% --------------------------------------------------------------------------
\subsection{Phase 5: Chain Verification --- Termination at Human}
\label{sec:rs-phase5-chain-verification}

The delegation chain is verified at any runtime point by the Chain-Verify algorithm (Algorithm~\ref{alg:chain-verify}), which confirms the structural validity of the chain from any node to its human anchor. This phase is not a post-hoc audit conducted after operations complete --- it is a runtime-verifiable structural property, checkable at any time by any actor with ledger read access.

\medskip
\noindent\textbf{Algorithm: Chain-Traverse.} The chain traversal algorithm builds the complete delegation lineage from any node to the human anchor by recursively applying $\alpha$:

\begin{algorithm}
\caption{Chain-Traverse($n$) --- Build the delegation chain from node $n$ to the human anchor}
\label{alg:chain-traverse}
\begin{algorithmic}[1]
\REQUIRE $n$ is a provisioned RSP node with $\alpha$ defined
\ENSURE $\Chain{n}$ --- the ordered set of all ancestor nodes including the human anchor
\STATE $chain \leftarrow \langle \rangle$ \COMMENT{Initialize empty ordered list}
\STATE $current \leftarrow n$ \COMMENT{Start at the query node}
\WHILE{$\text{true}$}
    \STATE $\text{append}(chain,\ current)$ \COMMENT{Add current node to chain}
    \IF{$\alpha(current) = \bot$}
        \STATE \textbf{break} \COMMENT{Root human anchor reached --- chain complete}
    \ENDIF
    \STATE $current \leftarrow \alpha(current)$ \COMMENT{Move to parent}
\ENDWHILE
\RETURN $chain$
\end{algorithmic}
\end{algorithm}

Chain-Traverse terminates in at most $D_{\max} + 1$ iterations because each application of $\alpha$ moves strictly closer to the root, and the root satisfies $\alpha(n) = \bot$. The output is uniquely determined for any input node because $\alpha$ is a total function with unique values.

\medskip
\noindent\textbf{Algorithm: Chain-Verify.} Chain-Verify confirms the structural validity of a delegation chain by checking four conditions:

\begin{algorithm}
\caption{Chain-Verify($n$) --- Verify the authorization chain from node $n$ to human anchor}
\label{alg:chain-verify}
\begin{algorithmic}[1]
\REQUIRE $n$ is a provisioned RSP node
\ENSURE $\text{Boolean}$ indicating whether the chain is a structurally valid RSP chain
\STATE $chain \leftarrow \text{Chain-Traverse}(n)$
\STATE $m \leftarrow \text{len}(chain)$

\STATE \COMMENT{\textit{Step V1: Verify chain terminates at human anchor}}
\IF{$\alpha(chain[m-1]) \neq \bot$}
    \STATE \RETURN $\text{false}$ \COMMENT{Root is not $\bot$ --- malformed chain}
\ENDIF
\IF{$\neg h(chain[m-1])$}
    \STATE \RETURN $\text{false}$ \COMMENT{Root is not a designated human principal}
\ENDIF

\STATE \COMMENT{\textit{Step V2: Verify each spawn event has a Purpose Layer warrant}}
\FOR{$i \leftarrow 0$ \TO $m-2$}
    \STATE $child \leftarrow chain[i]$
    \STATE $parent \leftarrow chain[i+1]$
    \IF{$\neg \text{warrant\_exists}(child,\ parent)$}
        \STATE \RETURN $\text{false}$ \COMMENT{Missing Purpose Layer warrant for spawn}
    \ENDIF
\ENDFOR

\STATE \COMMENT{\textit{Step V3: Verify scope refinement at each link}}
\FOR{$i \leftarrow 0$ \TO $m-2$}
    \STATE $child \leftarrow chain[i]$
    \STATE $parent \leftarrow chain[i+1]$
    \IF{$\neg (R(child) \subset R(parent))$}
        \STATE \RETURN $\text{false}$ \COMMENT{Scope refinement constraint violated}
    \ENDIF
\ENDFOR

\STATE \COMMENT{\textit{Step V4: Verify depth bound compliance at each node}}
\FOR{$i \leftarrow 0$ \TO $m-1$}
    \IF{$\text{depth}(chain[i]) > D_{\max}$}
        \STATE \RETURN $\text{false}$ \COMMENT{Depth bound exceeded at node}
    \ENDIF
\ENDFOR

\RETURN $\text{true}$
\end{algorithmic}
\end{algorithm}

Chain-Verify returns \texttt{true} if and only if all four verification conditions hold simultaneously. V1 confirms the chain terminates at a human anchor; V2 confirms every spawn was authorized by the Purpose Layer; V3 confirms scope refinement was respected at every link; V4 confirms the depth bound was respected throughout. The algorithm is deterministic and produces the same result regardless of which node initiates verification --- there is no self-serving interpretation of authorization that a drifted or orphaned node could produce.

\medskip
\noindent\textbf{Exclusive human terminus.} The chain terminates at the human anchor $h$ by architectural constraint. The human anchor is non-collapsing and non-substitutable. No machine actor can serve as the terminus of an RSP chain --- the chain must reach a Purpose Layer entity with a human principal designation, or it is not an RSP chain. This exclusivity is what makes the \bxthree{} upstream accountability guarantee structurally operative in recursive agent populations. The constant-time escalation bound provided by $D_{\max}$ ensures that reaching the human anchor from any node requires at most $D_{\max} + 1$ pointer dereferences, independent of the total number of active nodes in the system.

% --------------------------------------------------------------------------
\subsection{Depth Management: Max Depth, Spawn Refusal, and Convergence}
\label{sec:rs-depth-management}

The combination of $D_{\max}$, spawn refusal at the depth limit, and convergence criteria constitutes the complete depth management system of RSP. These are not three independent mechanisms --- they are three layers of a single depth enforcement architecture that ensures every RSP chain terminates within a finite number of steps, at or before the architectural depth bound, and always at a human principal.

\medskip
\noindent\textbf{Maximum spawn depth $D_{\max}$.} $D_{\max}$ is set by the Purpose Layer at system initialization based on the operational consequence domain of the deployment. High-stakes environments --- financial orchestration, clinical decision support, critical infrastructure --- require lower $D_{\max}$ values. The value is recorded in the Fact Layer authorization ledger and is immutable. The depth bound is not a suggestion --- it is a structural invariant enforced at every spawn step by both the Bounds Engine and the Fact Layer pre-commit hook.

\medskip
\noindent\textbf{Spawn refusal at limit.} When a spawn event would cause $\text{depth}(n_i) + 1 > D_{\max}$, the Bounds Engine refuses the spawn and sets the node to \texttt{ESCALATING} state. This is the spawn refusal condition. The refusal is deterministic: the Bounds Engine has no discretionary authority to exceed $D_{\max}$ regardless of operational urgency, task criticality, or the absence of an alternative resolution path. The refusal is logged as a ledger entry and the human anchor is notified through the escalation path. The chain does not proceed, does not bypass the depth bound, and does not continue autonomously past $D_{\max}$. The $D_{\text{ESCALATE}}$ threshold provides advance warning --- human confirmation is required before the depth limit is reached, not only after it is breached.

\medskip
\noindent\textbf{Convergence criteria.} A chain is said to have \emph{converged} when all of the following hold simultaneously:

\begin{enumerate}
    \item[(C1)] \textbf{Terminal state reached.} All leaf nodes are in \texttt{COMPLETED}, \texttt{TERMINATED}, or \texttt{ESCALATING} state with a \texttt{REJECT} directive from $h$.

    \item[(C2)] \textbf{No active exceptions.} No exception condition remains unresolved within the chain's operating scope.

    \item[(C3)] \textbf{Ledger continuity maintained.} All state transitions and spawn events are recorded in the Fact Layer forensic ledger with hash-chain continuity intact.

    \item[(C4)] \textbf{Verified human termination.} The chain is verified to terminate at the human anchor $h$ --- the root is a Purpose Layer entity with a human principal designation and $\alpha(n_0) = \bot$.
\end{enumerate}

Convergence is assessed by the Bounds Engine at each state transition event. When all four criteria hold, the chain is \emph{resolved}: no further autonomous spawn events are authorized, and the chain is archived as a completed ledger entry. The forensic record preserves the complete chain topology and all state transitions for post-hoc audit. Convergence is not merely the absence of active nodes --- it is a positive confirmation that the chain has reached a defined terminal state with a complete, tamper-evident accountability record that traces to a human principal.

\medskip
The Recursive Spawning Protocol, taken as a whole, provides a single architectural guarantee: that every agent in a multi-agent system operates within a delegation chain that is simultaneously immutable, auditable, scope-refined, depth-bounded, and verified to terminate at a human principal. The guarantee is structural --- it holds regardless of the goodness, reliability, or intentions of any machine actor in the chain. This is the operational expression of the upstream \bxthree{} accountability guarantee applied to recursive agent population dynamics.
